\documentclass[11pt]{article}
\usepackage[margin=1in]{geometry}
\usepackage{amsmath,amssymb}
\usepackage{booktabs}
\usepackage{hyperref}
\usepackage{microtype}
\usepackage{parskip}
\usepackage{xcolor}
\usepackage{listings}
\usepackage{graphicx}

\hypersetup{
  colorlinks=true,
  linkcolor=blue,
  citecolor=blue,
  urlcolor=blue,
}

\lstset{
  basicstyle=\ttfamily\small,
  breaklines=true,
  frame=single,
  backgroundcolor=\color{gray!8},
}

\title{HAMesh: Emergent Attractor Dynamics and Cross-Domain Synthesis\\
in Holographic Associative Memory Operating\\
in Large Language Model Embedding Space}

\author{Jesse Wallace\\
\textit{Independent Research}\\
\href{https://github.com/mescgit/HAMesh}{github.com/mescgit/HAMesh}}

\date{April 2026}

\begin{document}

\maketitle

\begin{abstract}
We introduce HAMesh, a holographic associative memory (HAM) that operates
natively in the embedding space of a large language model (LLM). Unlike
retrieval-augmented generation (RAG), which fetches discrete document chunks,
HAMesh \emph{synthesizes}: a query vector diffracts through the superposition
of all stored knowledge simultaneously, and the answer emerges through
constructive interference. Through recursive self-modification we call
\emph{dreaming}, the mesh develops stable attractor basins without explicit
supervision. We document four principal findings: \textbf{(1)}~attractor
identity is perfectly deterministic --- independent dream runs on the same
mesh converge to identical concept clusters (overlap = 1.000, $n=6$);
\textbf{(2)}~attractor clusters exhibit significantly higher semantic
coherence than random memory samples ($+0.22$ to $+0.30$ above baseline);
\textbf{(3)}~a multi-mesh collective achieves 100\% cross-domain activation in
curiosity-driven synthesis (16/16 insights draw from both specialised meshes
simultaneously); and \textbf{(4)}~a post-distillation phase transition
produces two deterministic equilibria --- an always-collapsed state at 0 dream
cycles and an always-distributed state at 500 cycles with mild decay --- with
stochastic bistability between, characteristic of systems near a bifurcation
point. HAMesh provides persistent, self-organising memory that grows through
use without retraining the underlying LLM.
\end{abstract}

% ============================================================
\section{Introduction}
% ============================================================

Large language models compress vast knowledge into fixed weights during
training. At inference time they are stateless --- each session begins with
no memory of prior interactions, and updating their knowledge requires
expensive retraining. Retrieval-augmented generation (RAG) partially addresses
this by maintaining an external vector store, but retrieval returns the
\emph{nearest stored chunk}, not a synthesis across the full knowledge base.

Holographic associative memories offer a different model of storage and
retrieval. In a HAM, knowledge is encoded as interference patterns in a dense
matrix through outer-product superposition. Retrieval is not search --- it is
diffraction: a query propagates through the matrix and the answer emerges from
the combined interference of all stored patterns simultaneously. This is
content-addressable, gracefully degrading, and fundamentally different from
lookup.

The key insight of HAMesh is to operate this holographic machinery \emph{inside
an LLM's native embedding space}. Rather than using random or learned
representation vectors, we use the LLM's own embedding function as the sensory
organ: queries are embedded by the LLM, diffracted through the holographic
mesh, and the synthesised output is decoded back to language by the same LLM.
The mesh inherits the LLM's full semantic geometry for free, enabling
meaningful interference between concepts the model understands.

We make four contributions:
\begin{enumerate}
  \item A holographic memory architecture that operates in LLM embedding space,
        with knowledge distilled from the LLM through generation rather than
        injected from external corpora.
  \item A recursive self-modification process (\emph{dreaming}) that develops
        stable attractor basins without supervision.
  \item A multi-mesh collective with cross-pollination and a curiosity engine
        for autonomous gap identification and cross-domain synthesis.
  \item An empirical characterisation of the phase transition between
        attractor-collapsed and multi-hop-capable mesh states, including two
        deterministic equilibria identified across six independent trials.
\end{enumerate}

% ============================================================
\section{Related Work}
% ============================================================

\paragraph{Holographic and associative memory.}
Kanerva's Sparse Distributed Memory~\cite{kanerva1988} and Plate's Holographic
Reduced Representations~\cite{plate1995} establish the mathematical foundations
for content-addressable memory via distributed superposition. HRR uses circular
convolution for compositional binding; our outer-product approach is closer to
Hopfield network storage~\cite{hopfield1982}.

\paragraph{Modern Hopfield networks.}
Ramsauer et al.~\cite{ramsauer2020} proved that transformer self-attention is
mathematically equivalent to dense associative memory retrieval, establishing a
formal connection between transformers and Hopfield networks. HAMesh can be
understood as an external Hopfield network operating in the same geometric space
as the transformer's attention, rather than inside it.

\paragraph{Holographic operations in LLM embedding space.}
Hypertokens / HDRAM~\cite{augeri2025} (2025) is the closest published work: it
treats the transformer's latent space as a spread-spectrum channel and applies
holographic operations for improved key-value retrieval ($2\times$ improvement).
HAMesh targets a different objective --- persistent synthesis and
self-organisation --- rather than retrieval efficiency.

\paragraph{Memory-augmented LLMs.}
MemGPT~\cite{packer2023}, LangMem, and Memori~\cite{borro2026} augment LLMs
with external memory, but all use retrieval-based approaches (nearest-neighbour
search over stored summaries). None develop self-organising topology or
attractor dynamics.

\paragraph{Dreaming in neural networks.}
Daydreaming Hopfield Networks~\cite{serricchio2025} explores dream-like updates
for correlated data. Prior work on sleep-inspired consolidation~\cite{tononi2003}
separates learning from unlearning phases. HAMesh extends dreaming to LLM
embedding spaces and documents the resulting attractor dynamics quantitatively.

% ============================================================
\section{Architecture}
% ============================================================

\subsection{The Holographic Mesh}

The core data structure is a $\mathit{dim} \times \mathit{dim}$ float32 matrix
$\mathbf{M}$, initialised to zero, where $\mathit{dim}$ is the LLM's embedding
dimension (4096 for Bonsai 8B). Knowledge is stored through \emph{folding}:
given a key embedding $\mathbf{k}$ and value embedding $\mathbf{v}$, both
L2-normalised:

\begin{equation}
  \mathbf{M} \leftarrow \mathbf{M} + s \cdot \hat{\mathbf{k}} \otimes \hat{\mathbf{v}}
\end{equation}

where $s$ is a strength scalar and $\otimes$ denotes outer product. Multiple
folds superimpose in the same physical space. This is Hebbian outer-product
storage --- mathematically equivalent to storing rank-1 updates to an
associative matrix.

\paragraph{Retrieval by diffraction.}
Given a query embedding $\mathbf{q}$:

\begin{equation}
  \text{output} = \text{normalize}(\mathbf{M}\,\hat{\mathbf{q}})
  = \text{normalize}\!\left(\sum_i s_i \cos(\mathbf{k}_i, \mathbf{q})\,\mathbf{v}_i\right)
\end{equation}

The output is a weighted combination of all stored value vectors, weighted by
cosine similarity between the query and each stored key. This is single-head
linear attention without softmax, operating holographically: \emph{all}
associations respond simultaneously. The result is not the stored value nearest
to the query --- it is a synthesis vector blended from every relevant pattern.

\paragraph{Multi-hop diffraction.}
Re-applying the diffraction operator to its own output enables transitive
reasoning:

\begin{align}
  \mathbf{h}_1 &= \text{normalize}(\mathbf{M}\,\hat{\mathbf{q}}) \quad (A \to B) \\
  \mathbf{h}_2 &= \text{normalize}(\mathbf{M}\,\mathbf{h}_1) \quad (B \to C)
\end{align}

Each hop follows associations one step further without explicit graph traversal.
We measure multi-hop \emph{diversity} as the fraction of queries where the
top-1 activated memory differs between 1-hop and 2-hop results.

\subsection{Knowledge Distillation Through Generation}

Rather than injecting external text, HAMesh extracts the LLM's own knowledge by
prompting it to generate explanations of seed topics, embedding both the topic
and explanation, and folding the association. For a seed topic $T$ with
generated explanation $E$:

\begin{equation}
  \mathbf{M} \leftarrow \text{fold}(\text{embed}(T),\; \text{embed}(E))
\end{equation}

This decompresses knowledge from the LLM's weights into holographic form. The
LLM generates subtopics, which are folded as children; cross-domain connections
suggested by the LLM are folded as bidirectional bridges at reduced strength.
We call this \emph{distillation through generation}.

\subsection{Dreaming}

The dreaming process implements recursive self-modification. Each cycle:
\begin{enumerate}
  \item Select a seed (random memory or previous cycle's output)
  \item Diffract through the mesh to get $\mathbf{u}$ (\emph{intuition})
  \item Fold: $\mathbf{M} \leftarrow \mathbf{M} + \alpha\,\mathbf{u} \otimes \mathbf{u}$ \quad (self-reference)
  \item Fold: $\mathbf{M} \leftarrow \mathbf{M} + \alpha\,\mathbf{s} \otimes \mathbf{u}$ \quad (bridge from seed)
  \item Cross-link top-2 activated memories at reduced strength
  \item Every 20 cycles apply decay: $\mathbf{M} \leftarrow (1-\delta)\,\mathbf{M}$
  \item Use $\mathbf{u}$ as next cycle's seed --- the loop
\end{enumerate}

Without decay ($\delta=0$), self-reference folds compound indefinitely. With
decay, the mesh reaches a dynamic equilibrium between self-reinforcement and
forgetting.

\subsection{Multi-Mesh Collective}

Two specialised meshes (science/technology and philosophy/history/nature) run
in parallel. Given a query embedding $\mathbf{q}$, each mesh diffracts
independently and their outputs are blended:

\begin{equation}
  \mathbf{b} = \text{normalize}\!\left(\sum_m \|\mathbf{M}_m\|_F \cdot \text{diffract}_m(\mathbf{q})\right)
\end{equation}

weighted by each mesh's Frobenius norm (more knowledge $\Rightarrow$ more
influence). Activated memories from both meshes are pooled and ranked by cosine
similarity to $\mathbf{b}$.

\textbf{Cross-pollination} transfers each mesh's dominant attractor patterns
(top eigenvector directions via power iteration) to the other at low strength
(0.04), creating cross-domain associations without homogenising the meshes.

\textbf{Curiosity engine.} Isolated memories (low average cosine similarity to
their $k$-nearest neighbours) represent knowledge islands. For each isolated
memory, the LLM generates probing questions at the knowledge boundary; the
collective diffracts each question (2-hop); the cross-domain answer is folded
into all meshes. This implements autonomous gap identification and filling.

% ============================================================
\section{Experiments}
% ============================================================

We evaluate five claims using Bonsai 8B (a 1-bit quantised Qwen3-8B, 1.15 GB
GGUF) running via llama.cpp on an RTX 4090. Two meshes were distilled:
science/technology (10 seed topics, depth 2, 5 subtopics each, $\sim$120
concepts) and philosophy/history/nature (same structure).

% ----------------------------------------
\subsection{C1: Attractor Identity Is Deterministic}

\textbf{Claim.} Independent dream runs on the same mesh converge to identical
attractor clusters.

\textbf{Method.} Three independent runs of 500 dream cycles ($\delta=0.02$)
on each mesh, loaded fresh from disk for each run. Measure pairwise overlap of
top-5 activated memories across run pairs.

\textbf{Results.}

\begin{table}[h]
\centering
\begin{tabular}{lccc}
\toprule
Mesh & Avg Pairwise Overlap & Avg Coherence & Verdict \\
\midrule
Science    & 1.000 & 0.733 & SUPPORTS \\
Philosophy & 1.000 & 0.667 & SUPPORTS \\
\bottomrule
\end{tabular}
\caption{C1 results: perfect attractor overlap across all 6 run pairs.}
\end{table}

Perfect overlap across all 6 run pairs. The attractors that formed --- a
cryptography/number-theory cluster in the science mesh and a
consciousness/cognition cluster in the philosophy mesh --- are identical
regardless of random seed order in dreaming.

% ----------------------------------------
\subsection{C2: Attractor Clusters Are Semantically Coherent}

\textbf{Claim.} The attractor clusters that form are meaningful --- not
high-energy noise, but real conceptual neighbourhoods.

\textbf{Method.} After dreaming to convergence, score the semantic coherence
of the top-5 attractor concepts using the LLM as judge (0--10 scale,
normalised). Compare against 10 random samples of 5 memories from the same
mesh.

\textbf{Results.}

\begin{table}[h]
\centering
\begin{tabular}{lcccc}
\toprule
Mesh & Attractor Coherence & Random Baseline & $\Delta$ & Verdict \\
\midrule
Science    & 0.80 & 0.50 & +0.30 & SUPPORTS \\
Philosophy & 0.70 & 0.48 & +0.22 & SUPPORTS \\
\bottomrule
\end{tabular}
\caption{C2 results: attractor coherence significantly exceeds random baseline.}
\end{table}

The science mesh's cryptography cluster (Key Exchange $\to$ Hashing $\to$
Number Theory $\to$ Computer Architecture) received coherence 0.80 --- number
theory is the mathematical foundation of cryptographic key exchange, a real
semantic relationship the mesh discovered through dreaming alone.

% ----------------------------------------
\subsection{C3: Cross-Pollination Creates Cross-Domain Curiosity}

\textbf{Claim.} After cross-pollination, each mesh develops curiosity about
topics from the other mesh's domain.

\textbf{Method.} After cross-pollination, we scan every memory in each mesh for
\texttt{[from X]} markers indicating foreign origin, then measure each such
memory's isolation score (mean cosine similarity to its 3 nearest neighbours).
Memories with isolation score below 0.3 are curiosity gaps; those above 0.3
are integrated bridges. This directly tests whether transferred knowledge sits
unintegrated (would activate curiosity) or merges smoothly into the host
topology.

\textbf{Results.} Cross-pollination transferred 16 cross-domain memories into
each mesh (32 total). All 32 had high connectivity scores (science mesh:
0.826--0.861; philosophy mesh: 0.959--0.969), placing them firmly in the
integrated category with zero isolation scores below 0.3. Contrary to the
initial hypothesis, dominant memories selected by power iteration are by
construction already well-connected in the source mesh and arrive pre-adapted to
connect in the target.

\textbf{Revision.} The original claim was too strong: cross-pollination does not
reliably create curiosity \emph{gaps}. Instead it creates cross-domain
\emph{bridges}. The curiosity engine finds gaps from \emph{distillation
imbalance} --- topics that were mentioned but never deeply folded --- rather
than from cross-pollination per se. We revise the claim to: \emph{cross-pollination
creates semantic bridges between domains; curiosity gaps emerge from the
asymmetry between what is known and what was only mentioned}.

\textbf{Verdict.} Partially supports C3 --- cross-domain memories are created
but integrate immediately rather than persisting as isolation targets.

% ----------------------------------------
\subsection{C4: Curiosity Synthesis Is Cross-Domain}

\textbf{Claim.} The curiosity engine generates insights drawing from multiple
mesh domains simultaneously.

\textbf{Method.} Run curiosity engine on 8 knowledge gaps; for each generated
Q$\to$A pair, record which meshes' memories activated above the confidence
floor.

\textbf{Results.}

\begin{table}[h]
\centering
\begin{tabular}{lc}
\toprule
Metric & Value \\
\midrule
Total insights & 16 \\
Cross-domain (both meshes activated) & 16/16 (100\%) \\
Integrity (real questions + answers) & 16/16 (100\%) \\
\bottomrule
\end{tabular}
\caption{C4 results: every insight activated both mesh domains.}
\end{table}

Every insight activated both the science and philosophy meshes. Sample
cross-domain syntheses:

\begin{itemize}
  \item \textit{``The non-commutative nature of matrix multiplication --- where
        order of multiplication affects the result --- mirrors the asymmetric
        nature of power relationships in philosophical and political discourse.''}
        [Gap: Matrix multiplication; Meshes: science $\times$ philosophy]

  \item \textit{``Eigenvalues offer an analogy to epistemological fixed points ---
        stable concepts that resist transformation under the application of new
        information.''}
        [Gap: Fermat's Last Theorem; Meshes: science $\times$ philosophy]

  \item \textit{``The constant motion of tectonic plates challenges traditional
        notions of stability and permanence --- a theme that resonates across
        existentialist and process philosophy traditions.''}
        [Gap: Plate tectonics; Meshes: science $\times$ philosophy]
\end{itemize}

% ----------------------------------------
\subsection{C5: Phase Transition in Multi-Hop Capability}

\textbf{Claim.} Multi-hop diversity undergoes a phase transition as a function
of dream cycles.

\textbf{Method.} Load fresh mesh, apply increasing dream cycles with
$\delta=0.02$, measure multi-hop diversity (fraction of 10 test queries where
2-hop top-1 differs from 1-hop top-1). Six independent trials.

\textbf{Results.}

\begin{table}[h]
\centering
\begin{tabular}{lcccc}
\toprule
Dream Cycles & Mean Diversity & Std Dev & $n$ & State \\
\midrule
0   & 0.0\%  & 0.0\%  & 6 & \textbf{Deterministic collapsed} \\
100 & 53.3\% & 48.9\% & 6 & Stochastic bistable \\
200 & 40.0\% & 42.4\% & 5 & Stochastic bistable \\
300 & 70.0\% & 44.7\% & 5 & Stochastic (mostly distributed) \\
400 & 60.0\% & 54.8\% & 5 & Stochastic bistable \\
500 & \textbf{100.0\%} & \textbf{0.0\%} & \textbf{6} & \textbf{Deterministic distributed} \\
700 & 60.0\% & 54.8\% & 5 & Stochastic bistable \\
1000 & 66.7\% & 51.6\% & 6 & Stochastic (mostly distributed) \\
\bottomrule
\end{tabular}
\caption{C5 phase transition results. Only 0 and 500 cycles are deterministic ($\sigma=0$, $n=6$).}
\end{table}

Two points are perfectly consistent across all trials: 0 cycles (always
collapsed) and 500 cycles (always distributed). All intermediate stages exhibit
high variance ($\sigma = 43$--$55\%$), indicating stochastic bistability --- the
random path taken through dream seed selection determines which attractor basin
is occupied. This is characteristic of systems near a bifurcation point, where
small perturbations can flip the system between competing stable states.

\textbf{Post-distillation collapse.} The fresh mesh (0 cycles) is already
collapsed despite no dreaming. During distillation, hub topics (physics,
computer architecture) receive disproportionate fold density --- each appears as
a seed, as parent of 5 subtopics, and as target in cross-links, accumulating
$\sim10\times$ the folds of peripheral topics. This creates attractor
concentration before dreaming begins. Dreaming with decay redistributes this
energy.

% ============================================================
\section{Discussion}
% ============================================================

\subsection{Synthesis vs.\ Retrieval}

RAG returns the stored chunk nearest to the query. HAMesh returns a synthesis
vector that is a weighted interference of all relevant stored patterns. The
practical difference is that HAMesh can respond to queries that fall
\emph{between} stored concepts --- the interference of ``genetics'' and
``computer architecture'' produces a response drawing from both, weighted by
relevance. No stored chunk captures this combination; it emerges from the
geometry of the superposition.

\subsection{The Goldilocks Zone}

The phase transition data (Table~5) has a practical implication: operate the
mesh at exactly 500 dream cycles with $\delta=0.02$ after distillation. This is
the only reliably distributed operating point. Fewer cycles leave the mesh in
the post-distillation collapsed state; more cycles re-introduce stochasticity
(though most paths remain distributed above 500 cycles). We recommend: distill,
dream exactly 500 cycles, save, then limit interactive sessions to $\le 100$
additional dream cycles.

\subsection{Biological Analogy}

The decay mechanism in dreaming is functionally analogous to synaptic
homeostasis during sleep~\cite{tononi2014}: the brain downscales synaptic
weights overnight to prevent runaway potentiation while preserving relative
connection strengths. HAMesh's decay prevents dominant attractors from
monopolising the mesh energy, maintaining the distributed state that enables
multi-hop reasoning. The Goldilocks zone may correspond to the point where
homeostatic decay and Hebbian strengthening reach equilibrium.

\subsection{Limitations}

\textbf{Capacity wall.} A $4096 \times 4096$ mesh can store $\sim$573
orthogonal patterns before interference noise degrades recall (Hopfield capacity
$\approx 0.14 \times \mathit{dim}$). With 250--370 stored memories, we are
approaching but not at this limit. Richer distillation may hit the wall.

\textbf{LLM dependence.} Without the LLM's embeddings and generation, the mesh
is a matrix of floats. The quality of synthesis depends entirely on Bonsai's
embedding geometry.

\textbf{C2 measurement variance.} Semantic coherence scoring via LLM judge
introduces variance; the $+0.22$--$+0.30$ delta is consistent but the absolute
scores vary run-to-run. A larger evaluation with diverse judges would strengthen
this claim.

\textbf{C5 bistability.} The high variance at intermediate cycle counts
($\sigma \approx 43$--$55\%$) means individual dream runs are not predictable
outcomes. The two deterministic anchor points are robust; the transition path is
not.

% ============================================================
\section{Conclusion}
% ============================================================

HAMesh demonstrates that holographic associative memory operating in LLM
embedding space exhibits rich emergent dynamics: deterministic attractor
formation, semantically coherent concept clustering, stochastic bistability
near a phase transition, and 100\% cross-domain synthesis in multi-mesh
curiosity-driven learning. The 500-cycle equilibrium point provides a practical
operating recommendation. The cross-domain synthesis results --- where a
curiosity engine operating on knowledge gaps in a science mesh autonomously
generates connections to philosophical fixed points and political asymmetry ---
suggest that multi-mesh architectures can produce genuinely novel cross-domain
reasoning without explicit instruction.

Code and experiment logs: \url{https://github.com/mescgit/HAMesh}

% ============================================================
\bibliographystyle{plain}
\begin{thebibliography}{10}

\bibitem{kanerva1988}
Kanerva, P. (1988).
\textit{Sparse Distributed Memory}.
MIT Press.

\bibitem{plate1995}
Plate, T.~A. (1995).
Holographic reduced representations.
\textit{IEEE Transactions on Neural Networks}, 6(3), 623--641.

\bibitem{hopfield1982}
Hopfield, J.~J. (1982).
Neural networks and physical systems with emergent collective computational
abilities.
\textit{Proceedings of the National Academy of Sciences}, 79(8), 2554--2558.

\bibitem{ramsauer2020}
Ramsauer, H., et al. (2020).
Hopfield Networks is All You Need.
\textit{arXiv:2008.02217}.

\bibitem{augeri2025}
Augeri, C.~J. (2025).
Hypertokens: Holographic Associative Memory in Tokenized LLMs.
\textit{arXiv:2507.00002}.

\bibitem{packer2023}
Packer, C., et al. (2023).
MemGPT: Towards LLMs as Operating Systems.
\textit{arXiv:2310.08560}.

\bibitem{borro2026}
Borro, L.~C., Macarini, L.~A.~B., Tindall, G., Montero, M., \& Struck, A.~B.
(2026).
Memori: A Persistent Memory Layer for Efficient, Context-Aware LLM Agents.
\textit{arXiv:2603.19935}.

\bibitem{serricchio2025}
Serricchio, L., Bocchi, D., Chilin, C., Marino, R., Negri, M., Cammarota, C.,
\& Ricci-Tersenghi, F. (2025).
Daydreaming Hopfield Networks and their surprising effectiveness on correlated
data.
\textit{Neural Networks}. arXiv:2405.08777.

\bibitem{tononi2003}
Tononi, G., \& Cirelli, C. (2003).
Sleep and synaptic homeostasis: a hypothesis.
\textit{Brain Research Bulletin}, 62(2), 143--150.

\bibitem{tononi2014}
Tononi, G., \& Cirelli, C. (2014).
Sleep and the price of plasticity: from synaptic and cellular homeostasis to
memory consolidation and integration.
\textit{Neuron}, 81(1), 12--34.

\end{thebibliography}

\end{document}
